\section{Evaluation Plan and Evidence Status}
\subsection{Available Model and Missing Mechanisms}
The existing platform provides capacity accounting, representative
inference DAGs, AXI packetization, finite MPU memory-service concurrency,
Ethernet credit flow control, shared HBM service, and a transaction-level
PCIe comparison path. These are behavioral transaction and packet models;
picosecond timestamp resolution does not make them instruction-level or
cycle-accurate GPU simulation.

Several gaps matter directly to the proposed claims. The current MPU
slot is released after memory service. Pending admission and pending
responses do not yet constitute a demonstrated complete finite-buffer
feedback chain. SUE data is accumulated at message granularity before
HBM writing, rather than using a complete packet-level credit/consumer
coupling. Aggregate issue limits must be defined per port or die, since
a per-task 1-ns issue interval does not establish a hardware issue rate.
Packet-rate counters do not impose a switch packet-processing limit.
If packet throughput is central to a result, the switch model needs
separate pipeline latency, initiation interval, parallel capacity, and
shared arbitration.

Current traces use synthesized routing and addresses, and do not
establish production locality. Object-level HBM residency, online
migration, continuous batching, and complete attention execution are not
assumed to be available. The proposed mechanism is not implemented in
this manuscript project.

\subsection{Calibration and Mechanism Experiments}
First, calibrate the model using measurements from real endpoints or
FPGA/RTL prototypes. Sweep request size, concurrency, read/write mix,
and response consumption rate, recording latency and throughput.
Separate fitting data from held-out verification cases. Measure both
service-slot occupancy and end-to-end occupancy, checking payload and
wire-byte conservation. If only individual subsystems can be calibrated,
state the resulting limit on system-level confidence.

Next, hold capacity and bandwidth fixed while varying service slots,
endpoint windows, receiver buffers, locality, and shared MoE traffic.
Access classes include active KV, session restoration, and fine-grained
lookup where supported by validated traces. This stage asks which
resource limits serviceable capacity, without claiming production
inference throughput.

\subsection{Baselines and Resource Fairness}
Baselines should include HBM-only with feasible batches, Host memory via
PCIe, all-remote MPU placement, uniform sharding, capacity-first,
bandwidth-proportional, and locality-based placement. Compare these with
\mechanism\ and small-instance static optima. Preserve accelerator count,
per-die HBM, remote capacity, links, service bandwidth, total slots,
buffers, and credits across policy comparisons.

Node-count scaling is a separate hardware-expansion experiment with an
explicit resource table. An HBM-only configuration that cannot admit a
large batch is a capacity rejection, not an infinite-latency performance
point. It should also be allowed to process smaller batches under the
same external workload. Cost or power claims require their own validated
models and are outside the initial scope.

\subsection{Ablations and Decisive Figures}
Model ablations progressively add memory-service concurrency, endpoint
windows, finite response storage, and shared resource timing to a
capacity/bandwidth baseline. Policy ablations remove placement or
transaction admission independently. Use unseen access patterns and
load levels after selecting parameters.

The primary proposed figure compares the physical-capacity frontier,
bandwidth-only prediction, detailed-model feasible region, and the
region attained by \mechanism. A second figure tests whether model
detail reverses a placement decision under identical hardware. Such a
reversal is a hypothesis, not a required outcome; numerical rounding or
unequal resources must not manufacture it. A third figure traces service
slots, response buffers, and endpoint budgets over time to explain any
observed boundary.

\subsection{Reporting and Falsification}
Initially report stage completion, transaction latency, useful and wire
bytes, occupancy, actual waiting time, and admitted working set. Summed
per-request waiting time is not a critical-path duration. Capacity
rejections, timeouts, and incomplete runs remain distinct from valid
performance points. End-to-end request goodput requires complete
inference dependencies, attention reads/writes, arrivals, and batching;
multiplying a representative block by layer count is insufficient.

The central architectural claim should be reduced or abandoned if
held-out workloads are adequately explained by bandwidth alone, if simple
baselines match the proposed policy, or if realistic finite buffering
eliminates the expected advantage. This draft provides no performance
results, new simulator executions, or hardware calibration evidence.
